\TNchapter[A certificate is not a passport; it is a lease. The agent's first launch is not its last. This chapter closes the book on the operational side of certification: documentation that stays current, version control that ties every change to a certificate, the triggers that force a recertification (model upgrade, material change, time elapsed), and the public disclosure obligations an enterprise will increasingly owe. It is also the chapter that argues --- quietly --- that the work of keeping an agent certified is most of the work of running a responsible agent program, and that the binder is also the playbook.]{Keeping an Agent Certified}
\label{ch:50-certification-20-keeping-certified}

\starttwocol

\section{Documentation as a living artifact}

The first thing your team will get wrong about certification, the day after launch, is to think it is finished. The binder closes. The owner moves to the next project. The model card, the SBOM, the behavior contract --- all of them age in place. Six months later, the agent your audit committee certified bears only a passing resemblance to the agent that is running in production.

The discipline is to treat every artifact in the binder as a living document. The model card has a \textit{last reviewed} date. The SBOM regenerates on every deploy. The behavior contract has a version number. The risk assessment is reopened on every material change. If your audit committee asks for the current state of any artifact and your team has to ``go check,'' the binder is stale.

The good news is that most of the artifacts can be wired to the systems that produce them. A model card can be generated from the model registry and the eval pipeline. An SBOM can be generated from the dependency graph. The traces and the evals can be queried from the observability stack. The work is in the wiring, not in the writing. Once wired, the binder refreshes itself.

\section{Version control for the certified system}

A certificate attaches to a specific version of an agent. The version is not just the model. It is the model, the prompts, the tools, the retrieval indices, the guardrails, and the dependencies --- the whole system whose SBOM you wrote in \autoref{ch:50-certification-18-pre-cert}. When any of these changes, the certified system is no longer the deployed system, and the binder no longer matches what is in production.

This is the place where most enterprises break their certification discipline first. The model vendor pushes a quarterly update. The product team tweaks the system prompt. An MCP server adds a tool. Each change is small. None of them, on its own, looks like a recertification event. Together, they have made the certificate obsolete.

The pattern that works is to tie the certificate to a \textit{configuration snapshot} --- a recorded, immutable record of every component the agent shipped with on the day it was certified --- and to require an explicit decision when any component changes. The decision is one of three: no recertification needed (the change is below the threshold); partial recertification (only the affected evidence needs to be refreshed); or full recertification (the change is material, the whole binder reopens). The threshold and the partial-recertification rules are themselves part of the binder, and they are reviewed annually.

\begin{TNwarning} Certifying an agent once and never again is the single most common failure mode of an enterprise certification program. The certificate that hangs on the wall a year after launch describes an agent that no longer exists. When the regulator or the incident review asks how the agent was approved for production, the answer ``we certified it in March'' is the wrong answer if it is now November and the model has been upgraded twice. \end{TNwarning}

\section{Recertification triggers}

The discipline of keeping a certificate current lives in \textbf{recertification triggers} --- the documented events that force the agent back through some or all of the certification process.

\begin{TNdefinition}{Recertification trigger} a documented condition (time elapsed, model upgraded, threshold breached, incident occurred) that forces an agent through certification again. Without triggers, certification is a one-time event; with them, it is a discipline. \end{TNdefinition}

The four trigger categories that cover most cases are summarized in the table below.

\begin{table*}[ht]
\centering
\caption{Recertification trigger categories.}
\small
\begin{tabularx}{\textwidth}{@{}l X l@{}}
\toprule
\textbf{Category} & \textbf{What fires it} & \textbf{Mechanism} \\
\midrule
\textbf{Time-based} & Twelve months elapsed (quarterly for high-tier); full binder refresh and review-board re-approval. & Calendar \\
\addlinespace
\textbf{Change-based} & Material change to the model, prompts, tools, data pipeline, or guardrails; ``material'' is defined in the binder. For EU AI Act systems: ``substantial modification'' per Article 43. Triggers partial recertification scoped to affected evidence. & Owner files \\
\addlinespace
\textbf{Metric-based} & A drop in a key production metric (accuracy, escalation rate, time-to-recover, hallucination rate) past a defined threshold. Fed by the continuous evaluation discipline. & Monitor \\
\addlinespace
\textbf{Incident-based} & A documented production incident, regardless of the calendar. The incident becomes part of the next binder. & Incident review \\
\bottomrule
\end{tabularx}
\end{table*}

The triggers are mechanical. The mechanism that turns them into action is a calendar, a monitor, and a named owner. The calendar fires the time-based triggers. The monitor fires the metric-based ones. The owner files the change-based ones. The board reviews them on a fixed cadence.

\begin{TNinpractice}{Aurora Retail} \textit{Aurora's dynamic-pricing agent had been certified in March and recertified in October, eighteen months later, after a model upgrade and a regulatory rule change.} The first recertification was painful --- six weeks, two thousand hours of staff time, three rounds of evidence assembly. The team treated it as a one-off. The second recertification was easier. The model card regenerated from the registry. The SBOM updated itself. The behavior contract had been versioned all along. The eval reports came out of the same pipeline that produced them in March. The cycle took two weeks and four hundred hours. The lesson was not that the second cycle was cheaper --- it was that the first cycle, done right, made the second one cheap. The binder had become the playbook. \end{TNinpractice}

For agents in scope of the EU AI Act, the post-market monitoring obligations of Articles 72 and 73 add a regulatory backstop to the metric-based and incident-based triggers. Article 72 requires providers of high-risk systems to operate a documented post-market monitoring system; Article 73 requires reporting of serious incidents within a defined window. Both feed your trigger calendar.

\section{Productisation: versioning, flags, and deprecation}

A certified agent is not a frozen artifact. It evolves --- new prompts, new tool versions, new client-specific overlays, new model vendors. The discipline of \textit{productising} an agent is what keeps the certificate honest as the product moves. Recertification triggers are the events that fire; productisation is the discipline that keeps the trigger fairly priced.

\textbf{Versioning.} Every certified agent has a version. \textit{Major} versions --- a new model, new core tools, a changed behavior contract --- trigger recertification. \textit{Minor} versions --- a prompt tweak, a threshold adjustment, a guardrail refinement --- trigger a lighter re-eval and a logged delta. \textit{Patch} versions --- bug fixes with no behavior change --- ship under the existing certificate. The discipline is to map every change to one of the three \textit{before} it ships, not after. Most certification programs fail this discipline by routing minor changes through the patch lane on the team's word, and discovering only at the next time-based trigger that the agent in production no longer matches the agent in the binder.

\textbf{Feature flags per BU or tenant.} Production agents serve more than one customer. Each BU or tenant gets its own feature-flag stack --- different escalation rules, different tool scopes, different retrieval indices, different refusal copy. Each flag stack is part of the certified configuration. A change to a flag's default is a version event, not a config tweak. The flags themselves live in the configuration snapshot the certificate attaches to, and the diff between two tenants' flag stacks is something the review board can actually read.

\textbf{Deprecation.} Every certified agent has a sunset. When the underlying model is end-of-life, when the business need shifts, when a successor agent supersedes it --- the deprecation procedure runs. It includes a notice window for tenants, a documented rollback plan, and a final post-mortem that closes the cert file rather than leaving it dangling on a shared drive. A certificate without a written end is a certificate that will outlive the agent it describes.

\begin{TNinpractice}{Cascade Health} \textit{Cascade's engineering-Q\&A agent ran for two years under the same certificate.} Over that time it accumulated sixteen feature-flag tweaks per BU, three minor versions, and one model-vendor migration. Each change passed the team's own version-mapping test on its way out the door. None tripped a recertification trigger. Until somebody noticed, eighteen months in, that the cert binder had not been re-read since the original sign-off. The post-mortem did not blame any individual change; it rewrote the productisation playbook --- a tighter definition of \textit{minor}, a per-tenant flag-stack diff in the binder, a quarterly read-through of the certificate by the agent's owner --- so the next cycle would not get there. The cycle after that one was clean. \end{TNinpractice}

The recertification triggers earlier in this chapter are what fire when versioning or deprecation lands; productisation is the discipline that decides what the trigger is actually pricing.

\section{Denial, conditions, and exceptions}

\autoref{ch:50-certification-19-process} named three outcomes the review board can return: approved, approved with conditions, denied. Recertification triggers fire on agents that are already certified. This section is about what happens when the outcome is not a clean approval, or when a recertification is required but cannot be performed.

\textbf{Denial.} A denial is not an indictment of the team. It is a finding that the evidence chain does not yet support the certificate. The remediation playbook is narrow: close the specific gap the board named, in writing, and resubmit. If the red-team report was missing, commission and run one. If the data lineage broke at a vendor boundary, document the boundary and the compensating control. The agent owner has an escalation path --- the audit committee or the standing certification authority --- if the team believes the board's finding is wrong, and that escalation is a formal request for review, not a re-litigation in the hallway. The resubmission timeline is tiered: thirty days for low-tier agents, sixty for medium, ninety for high. Meridian Logistics denied its dispatch agent in \autoref{ch:50-certification-19-process} for a trace-coverage gap; the team rebuilt the observability stack and resubmitted at the sixty-day mark. The denial was the cheap version of the incident the board did not want to read about later.

\textbf{Conditional certification.} \textit{Approved with conditions} is the outcome the previous chapter said teams underuse. In practice, conditions take three shapes. A \textit{time-bound deadline} --- conditional approval expires in ninety days unless the data-lineage gap is closed --- which puts the certificate on a clock the calendar enforces. A \textit{scope restriction} --- the agent ships, but only to a limited user population, or only for a subset of intents --- which lets the team learn in production without taking on the full risk surface. A \textit{mandatory monitoring interval} --- weekly drift report instead of monthly, with a named reviewer --- which raises the observability cost in exchange for a closer look. Conditions are written into the binder; the next review board reads them and confirms each one was met before lifting the conditional flag.

\textbf{Exception handling.} Sometimes a recertification trigger fires and the certification cannot be performed. The vendor model the binder depended on has been deprecated and the replacement is not yet on the approved list. The regulator has signaled new implementation guidance is coming and the team would rather wait than recertify against the old guidance. The evidence chain depends on a system that is offline for migration. The principle is the same in every case: the certificate \textit{freezes}. The system can no longer be verified, so the certificate stops asserting what it cannot back up. A freeze is paired with two things --- a documented rollback plan and a live kill switch (\autoref{ch:20-hardening-08-production}) --- and the freeze itself is logged, dated, and owned. The book's stance is plain: a frozen certificate is not a working certificate, and a frozen agent is not a certified agent. It is an agent on probation.

\begin{TNwarning} A long-lived frozen certificate is technical debt with regulatory exposure. The freeze that was a reasonable response to a vendor outage in week one becomes an unanswered question for the audit committee by month six. Set a maximum freeze window in the binder, in advance, before the first incident forces one. When the window expires, the agent either recertifies, ships under documented exception with explicit risk acceptance, or comes out of production. ``Indefinitely frozen'' is not a policy; it is the absence of one. \end{TNwarning}

\section{Public disclosure and the registry}

The last piece of the lifecycle is what the rest of the world sees. For low-tier agents, public disclosure is light --- a behavior contract on a marketing page, a privacy notice. For high-tier agents, especially those in regulated domains, disclosure is heavier and increasingly required.

Three artifacts carry the disclosure weight: a \textit{system card}, a \textit{transparency report}, and (for EU high-risk systems) a \textit{public registry entry}.

A \textbf{system card} is a model card's bigger cousin. Where the model card describes a model, the system card describes the deployed system --- the model, the tools, the populations served, the safeguards. It is the public version of much of what the binder contains. OpenAI, Anthropic, and other vendors now publish system cards for their major releases; large enterprises will increasingly publish them for their customer-facing agents.

A \textbf{transparency report} is a recurring (typically annual or quarterly) public summary of what the agent did and how it performed. The categories that belong in such a report are the transparency obligations introduced in the alignment chapters: what the agent did, what it refused, what it escalated, and what failed. The reader is your customer, your regulator, or your civil-society watcher. The goal is not perfect candor --- it is verifiable candor.

The \textbf{public registry entry} is the EU AI Act's specific obligation for Annex III high-risk systems. Article 71 of the Act requires registration in an EU-maintained public database before deployment, with the registration kept current through the agent's lifecycle. The data points the database collects mirror the model card and the system card. The discipline of keeping the registry entry accurate is the discipline of keeping the binder current.

\begin{figure*}[ht]
\centering
\resizebox{\textwidth}{!}{%
\begin{tikzpicture}[
  font=\sffamily\small,
  stage/.style={draw=TNNavy, line width=0.6pt, rounded corners=2pt,
    fill=TNNavyTint, text width=2.35cm, align=center, inner sep=7pt,
    minimum height=1.3cm, font=\sffamily\bfseries\color{TNNavy}},
  arr/.style={-{Stealth}, TNTealDeep, line width=0.7pt},
  arrlabel/.style={font=\sffamily\footnotesize\itshape\color{TNBodyText},
    align=center, text width=2.05cm}
]
  % ── Top row: forward certification path ──────────────────────────────
  \node[stage] (a) at ( 0.0,  0) {Develop};
  \node[stage] (b) at ( 4.1,  0) {Pre-cert};
  \node[stage] (c) at ( 8.2,  0) {Certify};
  \node[stage] (d) at (12.3,  0) {Operate};

  % ── Bottom row: Trigger and Re-certify sit below the forward path
  \node[stage] (e) at ( 4.1, -4.5) {Trigger};
  \node[stage] (f) at ( 8.2, -4.5) {Re-certify};

  % Top-row forward arrows; labels sit above the rectangles in open space.
  \draw[arr] (a.east) -- (b.west);
  \draw[arr] (b.east) -- (c.west);
  \draw[arr] (c.east) -- (d.west);
  \node[arrlabel] at ( 2.05, 0.95) {evidence};
  \node[arrlabel] at ( 6.15, 0.95) {review};
  \node[arrlabel] at (10.25, 0.95) {verdict};

  % Operate → Trigger: route below the top row before turning left
  \draw[arr]
    (d.south) -- ++(0,-1.2) -| (e.north);
  \node[arrlabel, text width=2.45cm] at (8.35, -1.55) {incident or drift};

  % Trigger → Re-certify: right along the bottom row (label below)
  \draw[arr] (e.east) -- (f.west);
  \node[arrlabel, text width=2.25cm] at (6.15, -5.45) {reassessment};

  % Re-certify → Operate: outer loop avoids the rectangles entirely
  \draw[arr, dashed]
    (f.south) -- ++(0,-0.7) -| (d.south);
  \node[arrlabel, text width=2.3cm] at (13.25, -3.45) {back to operate};
\end{tikzpicture}%
}
\caption{The certification lifecycle. Certification is not a one-time stamp; an incident or material change loops the agent back through reassessment.}
\end{figure*}

\begin{TNnew}
\section{Retiring the agent}

Deprecation gets one paragraph in most playbooks. It deserves more, because the agent you retire in year three may have two years of users who built workflows around it, a binder that the regulator still wants to read, and a data residue that does not go away when the endpoint does.

The cleanest public scaffold is the UK Government Service Manual's guidance on retiring a service \citep{gds-retiring-service}. It is short on AI-specific content and exactly right on the human posture. Three obligations carry over without modification.

\textbf{Tell users three things, on a timeline they can act on.} What is changing and why, what they have to do to continue to have the need met, and what happens to their data. Internal B2B users typically need 60 to 90 days. External API users need longer --- Nordic APIs' sunset writeups and the major-vendor playbooks (Salesforce, Mulesoft) cluster on 6 to 12 months \citep{nordic-apis-sunset}, and the GOV.UK PaaS decommission ran 18 months. The number is set by the worst-case integrator, not the median one.

\textbf{Keep the endpoints alive long enough that the transition is graceful.} The Service Manual asks for redirects to a successor for at least one year. The agent equivalent is a \textit{deprecation mode} --- the agent answers with a referral to the human or successor service, logs the call, and refuses the action --- for at least one full review-board cycle. The cost is small. The cost of cutting the cord on day one is the support queue you did not budget for.

\textbf{The binder outlives the agent.} The NIST AI RMF Playbook for the Govern function \citep{nist-ai-rmf-playbook-govern} calls for retention of ``ancillary data or artefacts such as predictions, explanations, intermediate input feature representations, and credentials'' after decommissioning. The EU AI Act is more specific: Article 18 requires the technical documentation be kept available to competent authorities \textbf{for ten years} after the system was placed on the market \citep{euAIact2024}. Article 12's logging obligation runs from deployment through decommissioning, not just through the current release. Decide before you retire which records go to cold storage, who can pull them, and what the lawful basis is for keeping the personal data inside them.

A retirement is a version event. It gets its own binder entry, its own review-board sign-off, and its own date on the public registry --- the same discipline the rest of the lifecycle taught.
\end{TNnew}

\section{Why the binder is also the playbook}

The binder is something a certifier reads. It is also something your team can run on. Done well, the eight artifacts of the pre-cert binder, the recertification triggers, the trace-and-eval discipline, and the review-board cadence are not extra work on top of running an agent. They are the work of running an agent --- observably, accountably, with the ability to say what it does and why.

This is the quiet argument the five pillars have been making all along. Alignment is what makes the agent do what you meant. Hardening is what bounds the cost of when it doesn't. Benchmarking is how you compare the agent to alternatives. Evaluation is how you measure it against the job you hired it for. Certification is how you defend it to the people who will hold you to account. The five pillars are not optional, they are not sequential, and they are not ``done'' --- they are continuously maintained capabilities. The first launch is the beginning, not the end.

Maya's customer-service agent --- the one with the \$4M budget request that opened this book --- does not need a 200-page certification binder before launch. It needs the eight artifacts, the right tier, the right standard, the right review board, the right triggers, and the discipline to keep all of them current. If your team can build and maintain that, you have a defensible agent program. If they cannot, you have a vendor problem.

The book closes by handing Maya, and you, a single sentence to take into the Monday meeting with the team: \textit{the binder is the playbook, the playbook is the program, and the program is what you are actually buying when you approve the project.}

\TNrule

\stoptwocol

\begin{TNkeytakeaways}
\item A certificate attaches to a specific version of an agent --- model, prompts, tools, data, guardrails. When any of those changes, the certified system is no longer the deployed system.
\item Recertification triggers (time, change, metric, incident) are what turn certification from a one-time event into a discipline.
\item The artifacts in the binder must be living documents, wired to the systems that produce them, not snapshots that age on a shared drive.
\item Public disclosure --- system cards, transparency reports, registry entries --- is becoming a baseline expectation, not a regulatory edge case.
\item The binder, well maintained, is also the playbook for running a responsible agent program; the work of certification and the work of operating are the same work.
\end{TNkeytakeaways}

\begin{TNchecklist}
\item Define the four trigger categories for each tier of agent and write them into the binder template.
\item Wire your model card and SBOM to regenerate on every deploy; treat manual updates as a failure mode.
\item Set a calendar for time-based recertification --- twelve months for medium tier, quarterly for high tier.
\item Tie one production metric per agent to a metric-based trigger; review the threshold every six months.
\item Decide, in writing, which model-vendor updates count as ``material'' for your portfolio.
\item Publish one transparency report this year --- quarterly is the goal, annual is the floor.
\item For any agent in EU AI Act Annex III scope, confirm the public-registry entry exists and is current.
\end{TNchecklist}



